% IOMED Article
% LaTeX Template — opendrg Whitepaper (IMRaD Structure)
%
% Build: tectonic wp_opendrg.tex
%
% License:
% CC BY-NC-SA 3.0 (http://creativecommons.org/licenses/by-nc-sa/3.0/)
%

%----------------------------------------------------------------------------------------
%	PACKAGES AND OTHER DOCUMENT CONFIGURATIONS
%----------------------------------------------------------------------------------------

\documentclass[fleqn,10pt]{IOMEDArticle}

% Fix for Tectonic/XeTeX fonts
\usepackage{ifxetex}
\ifxetex
  \usepackage{fontspec}
  \setmainfont{Times}
  \setsansfont{Helvetica}
\fi

\usepackage[english]{babel}
\usepackage{listings,color}
\usepackage{longtable}
\usepackage{array}
\usepackage{calc}
\usepackage{amsmath}
\usepackage{tikz}
\usetikzlibrary{arrows.meta,positioning,shapes.geometric,fit,backgrounds}
\hbadness=10000
\emergencystretch=1.5em
% Let long URLs in the reference list break anywhere.
\makeatletter
\g@addto@macro{\UrlBreaks}{\UrlOrds}
\makeatother

\lstset{
  basicstyle=\ttfamily\scriptsize,
  columns=fullflexible,
  keepspaces=true,
  breaklines=true,
  frame=single,
  framerule=0.3pt,
  rulecolor=\color{color1!60},
  backgroundcolor=\color{color2!4},
  xleftmargin=4pt,
  xrightmargin=4pt,
  aboveskip=4pt,
  belowskip=4pt
}

%----------------------------------------------------------------------------------------
%	ARTICLE INFORMATION
%----------------------------------------------------------------------------------------

\JournalInfo{IOMED Research \& Development}
\Archive{Whitepaper --- \today}
\PaperTitle{opendrg: Open, Portable Hospital Casemix Grouping in R}
\PaperSubtitle{Reproducing the CMS MS-DRG, NHS England HRG4+ and Prescribed Specialised Services classifications from official published rule tables}

\Authors{
    Gabriel Maeztu\textsuperscript{1,*}, Gianluca Fabiano\textsuperscript{2}, Rafael Pinedo-Villanueva\textsuperscript{2}
}

\affiliation{
    \textbf{1} IOMED Medical Solutions S.L., Barcelona, Spain \\
    \textbf{2} Health Economics and Outcomes Research (HEOR), Nuffield Department of Orthopaedics, Rheumatology and Musculoskeletal Sciences (NDORMS), University of Oxford, Oxford, UK (\url{https://github.com/HEOR-NDORMS}) \\
    $^*$ \textbf{Contact}: \texttt{gabriel.maeztu@iomed.health}; \texttt{gianluca.fabiano@ndorms.ox.ac.uk}; \texttt{rafael.pinedo@ndorms.ox.ac.uk}
}

\Keywords{Casemix --- Diagnosis Related Groups (MS-DRG) --- Healthcare Resource Groups (HRG4+) --- Prescribed Specialised Services --- Health Economics --- OMOP Common Data Model --- Open-Source Software}
\newcommand{\keywordname}{Keywords}

\Abstract{
\textbf{Background:} Hospital casemix groupers turn coded patient stays---diagnoses, procedures, demographics, and admission and discharge details---into standardised clinical and costing categories such as Medicare Severity Diagnosis Related Groups (MS-DRGs) in the United States and Healthcare Resource Groups (HRG4+) and Prescribed Specialised Services (PSS) flags in England. Because official groupers are distributed as closed, platform-restricted programs (Windows desktop applications, mainframe modules, or specialised Java libraries), they are difficult to embed in analytical workflows, incompatible with federated hospital databases where patient records cannot be exported, and opaque to researchers.
\textbf{Objective:} We developed \textbf{opendrg}, an open-source R package that reproduces MS-DRG, HRG4+, and PSS directly from official published rule tables, connects read-only to OMOP Common Data Model (CDM) databases, and evaluates cross-country portability on live European hospital data.
\textbf{Methods:} Official rule tables published by CMS and NHS England are imported without hand-coding clinical rules and evaluated in each agency's exact order. We verified MS-DRG against CMS's official software on all 9{,}864 \texttt{TESTDB} stays (v43.0, v43.1, v44.0), HRG4+ across all seven NHS care settings (2{,}250 official test records; 2024/25 National Costs and 2025/26 Local Payment), and PSS on all 500 official sample records. We further designed a two-country real-world OMOP portability study: completed across seven years (2018--2024; 35{,}312 inpatient stays) in a Spanish specialist hospital's OMOP CDM against official 3M APR-GRD cost weights and length of stay, with the parallel English confirmation on CPRD Aurum--HES OMOP (HRG4+ versus NHS England's Windows grouper) planned with the Oxford HEOR-NDORMS team.
\textbf{Results:} \textbf{opendrg} achieved \textbf{100\% record-by-record agreement} with CMS's official software across all three MS-DRG versions and \textbf{100\% agreement} with NHS England's expected outputs across all seven HRG4+ datasets and both PSS versions. On the live Spanish OMOP CDM, complete OMOP mapping grouped \textbf{99.8\%--100.0\% of stays every year} (53\%--59\% surgical), aligned 90\%--100\% of patients on major clinical categories with 3M APR-GRD, explained \textbf{72\%--81\% of out-of-sample cost-weight variation} (Spearman $\rho = 0.78\text{--}0.86$), and matched annual hospital Case-Mix Index within $\pm 0.2\%$; parallel UK CPRD--HES confirmation results are pending.
\textbf{Conclusion:} Treating published official rule tables as executable data delivers portable, auditable, 100\%-verified MS-DRG, HRG4+, and PSS grouping anywhere R runs, including directly inside federated OMOP hospital databases.
}

%----------------------------------------------------------------------------------------

\begin{document}

\flushbottom
\maketitle
\thispagestyle{empty}
\clearpage

\tableofcontents

%----------------------------------------------------------------------------------------
%	1. BACKGROUND
%----------------------------------------------------------------------------------------

\section{Background}

Diagnosis-related groups (DRGs) were introduced at Yale in the late 1970s to summarise hospital activity into a manageable number of clinically coherent, resource-homogeneous patient groups~\cite{fetter1980drg}. Four decades later, their descendants underpin hospital financing, performance benchmarking, risk adjustment, and health economics and outcomes research (HEOR) across most developed health systems.

Despite their central role, the official software tools that assign these classifications have remained closed, platform-restricted black boxes:
\begin{itemize}[leftmargin=*,nosep]
    \item \textbf{CMS MS-DRG} is distributed as mainframe software, a Windows desktop utility, and a Java library backed by proprietary binary files. Embedding it in standard statistical workflows such as R or Python requires custom technical wrappers or manual file exports.
    \item \textbf{NHS England HRG4+} is distributed exclusively as a Windows desktop program reading encrypted reference files. Costing and research teams working on macOS, Linux, or secure cloud research environments cannot run it at all.
    \item \textbf{Incompatible with federated research.} Both tools assume patient-level records are exported to a local file and fed into standalone software. In modern multi-centre research networks built on the OMOP Common Data Model~\cite{ohdsi2024omop} (OHDSI, DARWIN~EU, EHDEN), patient-level records never leave the hospital; analytical code travels to the database and only aggregate results return.
\end{itemize}

\subsection{United States: CMS MS-DRGs}

An MS-DRG is assigned to an acute inpatient stay from its principal diagnosis (main reason for admission), up to 24 secondary diagnoses with present-on-admission (POA) indicators, up to 25 procedures (ICD-10-CM and ICD-10-PCS), sex, and discharge status~\cite{cms2026msdrg,cms2026msdrgmanual}. Grouping screens first for resource-intensive \emph{pre-MDC} care (organ transplants, ECMO, tracheostomy with prolonged ventilation), assigns one of 25 body-system chapters (\emph{Major Diagnostic Categories}, MDCs) from the principal diagnosis, partitions stays into surgical versus medical classes, and splits by complication severity---with a Major Complication or Comorbidity (\emph{MCC}), with a Complication or Comorbidity (\emph{CC}), or without CC/MCC---after removing comorbidities redundant with the principal diagnosis. Unrelated operating-room procedures reroute to DRGs 981--989, and hospital-acquired conditions (HACs) not present on admission are stripped from severity in a second pass~\cite{cms2005hac}. Version 44.0 (FY 2027) defines 766 MS-DRGs plus two ungroupable codes (998 and 999).

\subsection{England: NHS England HRG4+ and Prescribed Specialised Services}

NHS England prices and costs activity across seven care settings---admitted patient care, outpatient (non-admitted) consultations, emergency medicine, adult, paediatric and neonatal critical care, and renal dialysis---through \textbf{HRG4+}~\cite{nhse2025hrg,nhsd2020casemix}. For admitted care, HRG4+ assigns a \emph{core} HRG to each \emph{finished consultant episode} (FCE, the period under one consultant) and to the overall hospital \emph{spell}, plus \emph{unbundled} HRGs for separately priced components (chemotherapy, radiotherapy, critical care, rehabilitation, specialist palliative care, complex imaging). Using ICD-10 diagnoses and OPCS-4 procedures, HRG4+ picks the dominant procedure from a national hierarchy, determines an HRG \emph{root}, and splits that root by complication-and-comorbidity (CC) score, age, or length of stay. The 2024/25 National Costs release contains 2{,}900 HRGs.

Alongside HRG4+, NHS England directly commissions over 150 \textbf{Prescribed Specialised Services} (rare cancers, transplantation, complex neurosurgery) and identifies qualifying episodes, outpatient attendances, and spells via published \emph{identification rules} (the PSO tool)~\cite{nhse2026pso}, assigning a dominant service line and its \emph{National Programme of Care} (NPoC).

\subsection{Spain: 3M/Solventum APR-GRD and CIE-10-ES}

Spain's National Health System (\emph{Sistema Nacional de Salud}, SNS) and regional health services (such as Catalonia's CatSalut) classify hospital discharges using \textbf{3M/Solventum's All Patient Refined DRGs (APR-GRD)} for acute inpatient stays---and International Refined DRGs (IR-GRD) for ambulatory surgery and day-hospital activity---run centrally by the Ministry of Health and regionally under commercial 3M Core Grouping Software licences~\cite{averill2003drg}. Extending the same DRG lineage as MS-DRG beyond Medicare's older adult population, APR-GRD adds finer neonatal and paediatric groups and attaches two four-level subclasses to every base DRG: \emph{Severity of Illness} (SOI, 1--4) for resource use and \emph{Risk of Mortality} (ROM, 1--4) for quality benchmarking.

Unlike MS-DRG and HRG4+, APR-GRD's grouping tables are proprietary and unpublished, so they cannot be openly re-implemented. Since 2016, however, Spanish hospitals have coded the national discharge register (RAE-CMBD, Royal Decree 69/2015) in \textbf{CIE-10-ES}---the official Spanish translation of US ICD-10-CM and ICD-10-PCS. Because the underlying clinical codes are identical, Spanish inpatient data can be grouped directly with an open MS-DRG engine inside the hospital's own database, making Spain an ideal real-world test of cross-country portability.

\subsection{Shared structure across published groupers}

Despite national differences in coding and terminology, MS-DRG, HRG4+, and PSO share the same underlying structure: \emph{ordered decision tables driven by code lists}. Each clinical code maps to clinical properties; ordered rules check which combination of properties a stay meets; and a severity or complexity score selects the final split. Because both CMS and NHS England already publish those tables openly, one transparent, data-driven architecture can reproduce all three faithfully.

%----------------------------------------------------------------------------------------
%	2. OBJECTIVE
%----------------------------------------------------------------------------------------

\section{Objective}

The goal of \textbf{opendrg}~\cite{opendrg2026} is to turn hospital casemix grouping from a platform-bound black box into a transparent, verifiable, and portable R~\cite{rcore2026} package (Figure~\ref{fig:architecture}, Table~\ref{tab:headline}) with three specific objectives:

\begin{enumerate}[leftmargin=*,nosep]
    \item \textbf{Reproduce official US and English classifications from published tables.} Implement CMS MS-DRG, NHS England HRG4+ (all seven care settings), and Prescribed Specialised Services (PSO) strictly from the agencies' published rule tables---with no hand-written clinical rules---and verify 100\% agreement against official software outputs.
    \item \textbf{Enable read-only federated grouping on OMOP CDM.} Provide unified OMOP adapters for both MS-DRG and HRG4+/PSO through \texttt{CDMConnector}~\cite{black2025cdmconnector} and \texttt{CohortConstructor}, keeping patient-level records inside the hospital database and reporting data-quality gaps explicitly.
    \item \textbf{Demonstrate real-world two-country portability on OMOP CDM.} Evaluate \textbf{opendrg} read-only on live European OMOP databases across two health systems: (a)~seven years of inpatient stays from a Spanish specialist hospital benchmarked against official 3M APR-GRD classifications, cost weights, and lengths of stay (completed), and (b)~linked English primary and hospital care (CPRD Aurum--HES OMOP) benchmarked against NHS England's official Windows HRG4+ Grouper and National Cost Collection unit costs (planned confirmation by Oxford HEOR-NDORMS).
\end{enumerate}

\begin{table*}[!t]\centering
\caption{Overview of the casemix classifications implemented and verified in \textbf{opendrg} (November 2026).}
\label{tab:headline}
\scriptsize
\setlength{\tabcolsep}{4.0pt}
\begin{tabular}{@{}p{2.2cm}p{2.5cm}p{3.6cm}p{4.8cm}p{3.0cm}@{}}
\toprule
Classification & Versions included & Source of official rules & Verification against official software & Status \\
\midrule
\textbf{CMS MS-DRG} (United States) & v43.0, v43.1 (FY 2026), v44.0 (FY 2027) & Nine published CMS rule tables + 14 hospital-acquired condition criteria & 9{,}864 / 9{,}864 CMS test stays identical to CMS's official software for every version: final \& initial DRG, MDC, med/surg, severity, return code & \textbf{Complete}; tested on Windows, macOS \& Linux \\
\textbf{NHS England HRG4+} (England) & 2024/25 National Costs, 2025/26 Local Payment & NHS England \emph{Code to Group} workbook (Open Government Licence v3.0) + published chapter flowcharts & 100\% agreement across all 7 care settings: 750 admitted episodes / 477 spells, 200 outpatients, 300 emergency, 800 critical care, 200 renal & \textbf{Complete} (episode \& spell HRGs, method, dominant procedure, unbundled HRGs, PBC, quality flags) \\
\textbf{NHS England Prescribed Specialised Services} (England) & PSO 2025/26, 2025/26 Version 2 & NHS England \emph{Guide to Identification Rules} workbook (OGL v3.0) & 100\% agreement on all official sample records: 300 inpatient episodes / 141 spells, 200 outpatient attendances, incl.\ validation messages & \textbf{Complete} (qualifying service lines, dominant line \& National Programme of Care) \\
\bottomrule
\end{tabular}
\end{table*}

\begin{figure*}[!t]\centering
\resizebox{\textwidth}{!}{%
\begin{tikzpicture}[
    font=\sffamily\scriptsize,
    box/.style={draw=color1!75, line width=0.7pt, rounded corners=3pt, fill=color2!6, align=center, inner xsep=0.2cm, inner ysep=0.14cm, text width=3.1cm, minimum height=1.25cm},
    data/.style={draw=color1, line width=0.9pt, rounded corners=3pt, fill=color1!10, align=center, inner xsep=0.2cm, inner ysep=0.14cm, text width=3.1cm, minimum height=1.25cm},
    res/.style={draw=color1, line width=0.9pt, rounded corners=3pt, fill=white, align=center, inner xsep=0.2cm, inner ysep=0.14cm, text width=3.1cm, minimum height=1.25cm},
    arr/.style={-{Stealth[length=2mm,width=1.6mm]}, line width=0.8pt, color=color1}
]
\node[data] (tables) at (0,0) {\textbf{Official rule tables}\\ Published CMS files \& NHS \emph{Code to Group}\\ imported once per release};
\node[box, right=0.55cm of tables] (mark) {\textbf{Code lookup}\\ Map diagnoses \& procedures to clinical categories, clusters \& exclusions};
\node[box, right=0.55cm of mark] (wide) {\textbf{Stay profile}\\ Principal \& secondary diagnoses, procedures, age, sex, discharge \& severity};
\node[box, right=0.55cm of wide] (eval) {\textbf{Ordered rule check}\\ Apply agency rules in exact official order (pre-MDC $\rightarrow$ MDC $\rightarrow$ splits)};
\node[res, right=0.55cm of eval] (res) {\textbf{Casemix output}\\ Final DRG / HRG, severity, unbundled groups, specialised flags \& audit};
\draw[arr] (tables) -- (mark);
\draw[arr] (mark) -- (wide);
\draw[arr] (wide) -- (eval);
\draw[arr] (eval) -- (res);
\node[below=0.35cm of mark, text width=3.1cm, align=center] {standard patient, diagnosis and procedure tables (or OMOP CDM)};
\node[below=0.35cm of eval, text width=3.1cm, align=center] {transparent tabular rules, no proprietary binaries};
\node[below=0.35cm of res, text width=3.1cm, align=center] {100\% verified against official CMS \& NHS England test data};
\end{tikzpicture}}
\caption{\textbf{Common architecture of \texttt{opendrg}.} All grouping rules reside in versioned official tables imported directly from CMS or NHS England. The engine looks up each stay's codes, builds a clinical profile, and evaluates the agency's decision rules in exact official order.}
\label{fig:architecture}
\end{figure*}

%----------------------------------------------------------------------------------------
%	3. MATERIALS AND METHODS
%----------------------------------------------------------------------------------------

\section{Materials and Methods}
\label{sec:methods}

\subsection{Published rule tables}

A core design principle of \textbf{opendrg} is separating \emph{rules} (versioned official data tables) from \emph{evaluation} (a small, unchanging engine). Updating to a new fiscal year requires importing the agency's newly published tables with one function call.

\textbf{CMS MS-DRG tables.} Alongside its software, CMS publishes nine plain-text tables defining the complete MS-DRG logic (Table~\ref{tab:flatfiles}, $\approx$170~KB compressed per version). Cross-checking these tables against CMS's binary release confirmed identical content and order across all 794 DRG rules, 8{,}271 procedure clusters, and 172{,}853 exclusion pairs. Only the 14 hospital-acquired-condition eligibility rules appear exclusively in CMS's binary files and are included as a documented per-version supplement.

\begin{table}[!t]\centering
\caption{Official CMS MS-DRG v44.0 rule tables bundled in \textbf{opendrg} ($\approx$170~KB compressed per version).}
\label{tab:flatfiles}
\scriptsize
\setlength{\tabcolsep}{3.5pt}
\begin{tabular}{@{}p{2.05cm}p{4.55cm}r@{}}
\toprule
Table & Clinical content & Rows \\
\midrule
\texttt{DRGLOGIC} & Ordered decision rules per DRG: chapter, DRG, surgical/medical class, reroutes, severity exclusions & 794 \\
\texttt{DXLIST} & All billable ICD-10-CM codes with sex limits and exclusion groups & 74{,}966 \\
\texttt{DXPATERN} & Clinical attributes, chapters and HAC lists per diagnosis group & 788 \\
\texttt{PRLIST} & All ICD-10-PCS procedure codes (and procedure clusters) & 83{,}380 \\
\texttt{PRPATERN} & Clinical attributes per procedure group & 1{,}042 \\
\texttt{DXATTLST}, \texttt{PRATTLST} & Named diagnosis and procedure clinical groups & 65 / 233 \\
\texttt{EXCLUSN} & Complication/comorbidity exclusion pairs (principal $\rightarrow$ excluded secondary) & 172{,}853 \\
\texttt{PRCLSTRS} & Multi-code procedure combinations (4{,}084 clusters) & 8{,}271 \\
\texttt{MDCDSC}, \texttt{DRGDSC4} & Official chapter and DRG titles & 25 / 1{,}034 \\
HAC rules & Eligibility criteria for the 14 hospital-acquired conditions & 14 \\
\bottomrule
\end{tabular}
\end{table}

\textbf{NHS England HRG4+ tables.} With every HRG4+ release, NHS England publishes a \emph{Code to Group} Excel workbook under the Open Government Licence~\cite{nhse2025hrg,ogl3}. Rather than mere prose documentation, the workbook contains the complete machine-readable rule tables powering the Windows grouper (Table~\ref{tab:hrgtables}, $\approx$272~KB compressed): ICD-10 and OPCS-4 mappings, clinical subchapters, combination procedures (primary operation plus site/approach qualifiers), the national procedure hierarchy, complication-and-comorbidity scores, programme budgeting categories, and formal condition flags governing every HRG root and split. The handful of decision trees published only as PDF flowcharts in the Chapter Summaries (emergency-medicine hierarchies and paediatric/neonatal critical-care trees) were transcribed into versioned reference tables bundled alongside the workbook tables.

\begin{table}[!t]\centering
\caption{Official HRG4+ 2024/25 National Costs rule tables imported from NHS England's \emph{Code to Group} workbook ($\approx$272~KB compressed).}
\label{tab:hrgtables}
\scriptsize
\setlength{\tabcolsep}{3.5pt}
\begin{tabular}{@{}p{2.9cm}p{3.9cm}r@{}}
\toprule
Table & Clinical content & Rows \\
\midrule
\texttt{hrgs} & HRG catalogue: 2{,}900 HRGs, 85 subchapters, 21 chapters (core \& unbundled) & 2{,}900 \\
\texttt{code\_to\_group} & ICD-10, OPCS-4 and combination codes & 39{,}474 \\
\texttt{code\_to\_group\_ triggers} & Clinical subchapters triggered per code & 56{,}404 \\
\texttt{code\_to\_group\_ hrgs} & Candidate HRG roots and conditions per code & 133{,}847 \\
\texttt{group\_to\_split\_ hrgs} & Complexity, age and length-of-stay splits per root & 2{,}944 \\
\texttt{flag\_rules} & Formal condition flags governing splits and roots & 1{,}788 \\
\texttt{lists} & Complication scores (158{,}466), procedure hierarchy (31{,}294), other lists & 247{,}614 \\
\texttt{combination\_codes} & Primary procedure + qualifying site/approach codes & 5{,}080 \\
\texttt{pbc\_mappings} & Programme budgeting categories & 28{,}816 \\
\bottomrule
\end{tabular}
\end{table}

\textbf{NHS England PSO tables.} NHS England's \emph{Guide to Identification Rules} workbook~\cite{nhse2026pso} enumerates all 214 specialised service lines, 261 identification rules (812 rule conditions over 629 code and eligibility lists), dominance hierarchies, and NPoC mappings; \textbf{opendrg} imports them into a single 42~KB table set.

\subsection{Grouping engines: MS-DRG and HRG4+ side by side}

Both engines follow the same transparent pipeline (Figure~\ref{fig:architecture}) while respecting each system's exact official rules.

\textbf{CMS MS-DRG workflow.} Every \texttt{DRGLOGIC} row expresses one candidate DRG as a yes/no condition over clinical groups, severity, and stay characteristics---for example, \emph{ENT malignancy with MCC} (DRG~146), \emph{ECMO or tracheostomy with prolonged mechanical ventilation} (DRG~3), or \emph{bilateral lower-limb joint replacement with MCC} (DRG~461). Grouping proceeds in six steps: (1)~map each ICD-10-CM diagnosis and ICD-10-PCS procedure to its CMS clinical groups, chapter, and sex-specific rules; (2)~fire multi-code procedure clusters and apply cluster suppression; (3)~remove secondary diagnoses excluded by the principal diagnosis or counted only if discharged alive; (4)~compute rule-specific severity (\emph{MCC}, \emph{CC}, or neither); (5)~evaluate rules top-to-bottom (first match wins) across pre-MDC, the principal diagnosis's MDC, and reroutes (multiple trauma, HIV, unrelated surgery); and (6)~run a second pass removing qualifying hospital-acquired conditions not present on admission to yield both \emph{initial} and \emph{final} DRGs.

\textbf{NHS England HRG4+ and PSO workflow.} For each admitted-patient-care episode, the HRG4+ engine: (1)~validates demographic, administrative, ICD-10, and OPCS-4 fields (flagging invalid episodes as \texttt{UZ01Z} with official quality messages); (2)~forms procedure combination codes from driving operations and their contiguous site/approach qualifiers; (3)~derives unbundled HRGs (consuming their driving procedures); (4)~selects the dominant procedure as the highest-ranked remaining code in NHS England's procedure hierarchy; (5)~evaluates candidate HRG roots under NHS precedence (burns $>$ multiple trauma $>$ procedure $>$ global exception $>$ diagnosis, last satisfied rule wins) and Core~3 escalations; (6)~applies complexity, age, and length-of-stay splits plus programme-budgeting mappings; and (7)~groups multi-episode \emph{spells} from the dominant episode combined with the union of the spell's diagnoses and procedures. The other six datasets (outpatients, emergency medicine, adult/paediatric/neonatal critical care, and renal dialysis) reuse the same rule engine with dataset-specific inputs and hierarchies. The PSO engine evaluates each service line's rule lists over the same episode or outpatient fields and picks the dominant qualifying service line and NPoC per episode, attendance, and spell.

Both systems share a uniform, three-line R interface:

\noindent\begin{minipage}{\linewidth}
\begin{lstlisting}[language=R]
library(opendrg)

# US / Spain (MS-DRG):
g <- grouper("msdrg", version = "44.0")
msdrg_out <- group(g, claims, diagnoses, procedures)

# England (HRG4+ admitted care, or dataset = "nac", "em", ...):
h <- grouper("hrg4plus", version = "2025-26-lp")
hrg_out <- group(h, episodes)
\end{lstlisting}
\end{minipage}

\subsection{Verification protocol against official software}

Written manuals describe clinical intent, whereas reimbursement and research reproducibility require matching the exact outputs of official software.
\begin{itemize}[leftmargin=*,nosep]
    \item \textbf{MS-DRG oracle:} We ran CMS's official Java grouper across CMS's complete 9{,}864-stay test database (\texttt{TESTDB}) for versions 43.0, 43.1, and 44.0 (first confirming that our setup reproduced CMS's own embedded expectations on all 9{,}864 stays), and compared every output field against \textbf{opendrg}, including HAC shifts and unknown hospital POA status.
    \item \textbf{HRG4+ and PSO oracles:} We compared \textbf{opendrg} record-by-record against NHS England's official expected outputs across all seven HRG4+ datasets (2{,}250 records in total, covering both the 2024/25 National Costs and 2025/26 Local Payment designs~\cite{nhse2025lpg}) and against NHS England's official PSO sample datasets (300 inpatient episodes / 141 spells and 200 outpatient attendances, plus verifying the exact May 2026 Version~2 rule differences).
\end{itemize}

\subsection{Grouping on hospital OMOP databases}

{\sloppy In federated research networks (OHDSI, DARWIN~EU, EHDEN), patient records remain inside each hospital's database and researchers connect via \texttt{CDMConnector}~\cite{black2025cdmconnector}. \textbf{opendrg} provides symmetrical OMOP adapters for both classifications---\texttt{omop\_msdrg()} for MS-DRG and \texttt{omop\_hrg()} for NHS England HRG4+ (plus direct PSO flagging via \texttt{group(grouper("pso"), cdm)})---all accepting either a full database reference or a study cohort created with \texttt{CohortConstructor}.\par}

Figure~\ref{fig:omop_er} shows the OMOP CDM tables read and the casemix outputs produced. All filtering and code lookups execute read-only inside the hospital database across both OMOP CDM v5.3 and v5.4. For MS-DRG, the adapter extracts ICD-10-CM and ICD-10-PCS codes and maps discharge status to UB-04 codes; for HRG4+ and PSO, it extracts ICD-10 (automatically padding three-character codes such as \texttt{I10} to NHS format \texttt{I10X}) and OPCS-4 (preserving recorded procedure order for combination codes), supports both one-episode-per-spell and multi-episode \texttt{visit\_detail} (FCE) grouping, and maps provider specialties and admission/discharge fields to NHS codes (\texttt{MAINSPEF}/\texttt{TRETSPEF}, \texttt{ADMIMETH}/\texttt{ADMISORC}, \texttt{DISMETH}/\texttt{DISDEST}).

\begin{figure*}[!t]\centering
\resizebox{\textwidth}{!}{%
\begin{tikzpicture}[
    font=\sffamily\scriptsize,
    ent/.style={draw=color1!85, line width=0.8pt, rounded corners=2.5pt, fill=white, align=left, inner xsep=0.18cm, inner ysep=0.12cm, text width=4.15cm},
    entopt/.style={draw=color1!65, dashed, line width=0.75pt, rounded corners=2.5pt, fill=color2!4, align=left, inner xsep=0.18cm, inner ysep=0.12cm, text width=4.15cm},
    anchorbox/.style={draw=color1, line width=1.15pt, rounded corners=2.5pt, fill=color1!8, align=left, inner xsep=0.18cm, inner ysep=0.13cm, text width=4.3cm},
    adapter/.style={draw=color1, line width=1.1pt, rounded corners=3.5pt, fill=color1!14, align=left, inner xsep=0.2cm, inner ysep=0.16cm, text width=3.85cm},
    outbox/.style={draw=color1, line width=0.95pt, rounded corners=2.5pt, fill=color2!8, align=left, inner xsep=0.18cm, inner ysep=0.12cm, text width=4.65cm},
    rel/.style={-{Stealth[length=1.9mm,width=1.4mm]}, line width=0.75pt, color=color1!85},
    relopt/.style={-{Stealth[length=1.9mm,width=1.4mm]}, dashed, line width=0.7pt, color=color1!65},
    flow/.style={-{Stealth[length=2.3mm,width=1.8mm]}, line width=1.05pt, color=color1},
    lbl/.style={font=\sffamily\tiny, fill=white, inner sep=1.2pt, text=color1!90}
]

% --- Column 2 (x = 6.2): Core Visits & Clinical Events ---
\node[ent] (cond) at (6.2, 4.6) {%
  \textbf{\texttt{condition\_occurrence}}\\[-1pt]
  \rule{\linewidth}{0.35pt}\\[-1pt]
  \texttt{PK condition\_occurrence\_id}\\
  \texttt{FK visit\_occurrence\_id, visit\_detail\_id}\\
  \texttt{FK condition\_source\_concept\_id}\\
  \texttt{FK condition\_concept\_id}\\
  \texttt{condition\_status\_concept\_id} (32902/03)\\
  \texttt{condition\_type\_concept\_id} (excl.\ 32858)%
};

\node[anchorbox, below=0.65cm of cond] (visit) {%
  \textbf{\texttt{visit\_occurrence}} \textit{(inpatient anchor)}\\[-1pt]
  \rule{\linewidth}{0.45pt}\\[-1pt]
  \texttt{PK visit\_occurrence\_id, FK person\_id}\\
  \texttt{visit\_concept\_id} $\in \{9201, 262\}$\\
  \texttt{visit\_start\_date, visit\_end\_date}\\
  \texttt{discharged\_to\_concept\_id / \_source\_value}\\
  \texttt{admitted\_from\_*} \textit{(HRG4+)}\\
  \texttt{FK provider\_id, care\_site\_id} \textit{(HRG4+)}%
};

\node[entopt, below=0.65cm of visit] (vdet) {%
  \textbf{\texttt{visit\_detail}} \textit{(optional HRG4+ FCEs)}\\[-1pt]
  \rule{\linewidth}{0.35pt}\\[-1pt]
  \texttt{PK visit\_detail\_id, FK visit\_occurrence\_id}\\
  \texttt{visit\_detail\_start/end\_date, FK provider\_id}%
};

\node[ent, below=0.65cm of vdet] (proc) {%
  \textbf{\texttt{procedure\_occurrence}}\\[-1pt]
  \rule{\linewidth}{0.35pt}\\[-1pt]
  \texttt{PK procedure\_occurrence\_id}\\
  \texttt{FK visit\_occurrence\_id, visit\_detail\_id}\\
  \texttt{FK procedure\_concept\_id, \_source\_concept\_id}\\
  \texttt{procedure\_date} \textit{(order kept for OPCS-4)}%
};

% --- Column 1 (x = 0): Patient, Cohort, Death, Admin ---
\node[entopt] (cohort) at (0, 4.6) {%
  \textbf{\texttt{cohort\_table}} \textit{(optional, CohortConstructor)}\\[-1pt]
  \rule{\linewidth}{0.35pt}\\[-1pt]
  \texttt{PK cohort\_definition\_id}\\
  \texttt{FK subject\_id} $\rightarrow$ \texttt{person}\\
  \texttt{cohort\_start\_date, cohort\_end\_date}\\
  \textit{attrs:} \texttt{cohort\_set}, \texttt{cohort\_codelist}%
};

\node[ent, below=0.75cm of cohort] (person) {%
  \textbf{\texttt{person}}\\[-1pt]
  \rule{\linewidth}{0.35pt}\\[-1pt]
  \texttt{PK person\_id}\\
  \texttt{gender\_concept\_id} (8507 / 8532)\\
  \texttt{birth\_datetime}\\
  \texttt{year / month / day\_of\_birth}%
};

\node[entopt, below=0.75cm of person] (death) {%
  \textbf{\texttt{death}} \textit{(optional)}\\[-1pt]
  \rule{\linewidth}{0.35pt}\\[-1pt]
  \texttt{PK/FK person\_id}\\
  \texttt{death\_date} $\rightarrow$ UB-04 \texttt{20} / \texttt{DISMETH=4}%
};

\node[entopt, below=0.75cm of death] (admin) {%
  \textbf{\texttt{provider} \& \texttt{care\_site}} \textit{(HRG4+/PSO)}\\[-1pt]
  \rule{\linewidth}{0.35pt}\\[-1pt]
  \texttt{PK provider\_id, specialty\_concept\_id}\\
  \texttt{specialty\_source\_value} $\rightarrow$ \texttt{MAINSPEF/TRETSPEF}\\
  \texttt{PK care\_site\_id, care\_site\_source\_value} $\rightarrow$ \texttt{PROCODET}%
};

% --- Column 3 (x = 12.2): OMOP Vocabularies (inside OMOP CDM box) ---
\node[ent] (concept) at (12.2, 2.2) {%
  \textbf{\texttt{concept}}\\[-1pt]
  \rule{\linewidth}{0.35pt}\\[-1pt]
  \texttt{PK concept\_id}\\
  \texttt{vocabulary\_id}:\\
  \quad\texttt{ICD10CM}, \texttt{ICD10PCS},\\
  \quad\texttt{ICD10}, \texttt{OPCS4},\\
  \quad\texttt{UB04 Pt dis status}\\
  \texttt{concept\_code}%
};

\node[entopt, below=1.0cm of concept] (crel) {%
  \textbf{\texttt{concept\_relationship}} \textit{(fallback)}\\[-1pt]
  \rule{\linewidth}{0.35pt}\\[-1pt]
  \texttt{FK concept\_id\_1, concept\_id\_2}\\
  \texttt{relationship\_id = "Maps to"}\\
  \textit{used when} \texttt{vocabulary\_fallback = TRUE}%
};

% --- Background grouping box strictly around OMOP CDM tables ---
\begin{pgfonlayer}{background}
  \node[draw=color1!45, dashed, rounded corners=5pt, fill=color2!2, inner xsep=0.38cm, inner ysep=0.35cm,
        fit=(cohort)(person)(death)(admin)(cond)(visit)(vdet)(proc)(concept)(crel)] (cdmbox) {};
  \node[anchor=south west, font=\sffamily\bfseries\scriptsize, text=color1] at ([yshift=2pt]cdmbox.north west)
        {OMOP CDM v5.3 / v5.4 Database (read-only via \texttt{CDMConnector})};
\end{pgfonlayer}

% --- Column 4 (x = 17.8): opendrg OMOP Adapter ---
\node[adapter] (adapter) at (17.8, 0.1) {%
  \textbf{opendrg OMOP Adapter}\\
  \texttt{omop\_to\_msdrg()}\\
  \texttt{omop\_to\_hrg()}\\[-1pt]
  \rule{\linewidth}{0.45pt}\\[-1pt]
  $\bullet$ Read-only database queries\\
  $\bullet$ Supports CDM v5.3 \& v5.4\\
  $\bullet$ Principal strategy:\\
  \quad\texttt{status}, \texttt{codelist},\\
  \quad\texttt{cohort\_codelist}, \texttt{all}, \texttt{first}\\
  $\bullet$ Code formatting \& \texttt{X}-padding\\
  $\bullet$ Discharge, specialty \&\\
  \quad readiness diagnostics%
};

% --- Column 5 (x = 25.0): Grouped Outputs ---
\node[outbox] (out_msdrg) at (25.0, 4.1) {%
  \textbf{MS-DRG Output} (\texttt{omop\_msdrg()})\\[-1pt]
  \rule{\linewidth}{0.4pt}\\[-1pt]
  \textit{1 row / stay (or cohort-stay / candidate):}\\
  \texttt{claim\_id, cohort\_definition\_id, cohort\_name}\\
  \texttt{visit\_occurrence\_id, person\_id, los, age, sex}\\
  \texttt{principal\_code, principal\_source}\\
  \texttt{drg, initial\_drg, base\_drg, mdc, med\_surg}\\
  \texttt{severity, return\_code, drg\_description}\\
  \texttt{n\_dx, n\_pr, has\_unmapped\_pr}%
};

\node[outbox, below=0.7cm of out_msdrg] (out_hrg) {%
  \textbf{HRG4+ / PSO Output} (\texttt{omop\_hrg()} / \texttt{group(pso)})\\[-1pt]
  \rule{\linewidth}{0.4pt}\\[-1pt]
  \textbf{\texttt{\$episodes}} \textit{(1 row / FCE or spell):}\\
  \texttt{episode\_id, cohort\_*, visit\_occurrence\_id,}\\
  \texttt{visit\_detail\_id, person\_id, fce\_hrg,}\\
  \texttt{grouping\_method, dominant\_procedure,}\\
  \texttt{unbundled\_hrgs, fce\_pbc, spell\_hrg}\\
  \textit{PSO:} \texttt{service\_line, fce\_npoc, spell\_service\_line}\\
  \textbf{\texttt{\$spells}}: \texttt{spell\_hrg, spell\_episode\_count, spell\_los}%
};

\node[ent, below=0.7cm of out_hrg] (out_qual) {%
  \textbf{Quality Report} (\texttt{omop\_*\_quality()})\\[-1pt]
  \rule{\linewidth}{0.35pt}\\[-1pt]
  Coding depth, primary-flag \%, billable vs\\
  headers / \texttt{X}-padded, unmapped procedures,\\
  discharge/specialty sources, grouping rate%
};

% --- ER Relationships ---
\draw[relopt] (cohort.east) -- ++(0.65,0) |- node[lbl, pos=0.22] {scope} ([yshift=0.75cm]visit.west);
\draw[rel] (person.east) -- ++(0.45,0) |- node[lbl, pos=0.22] {1 : *} ([yshift=0.2cm]visit.west);
\draw[relopt] (death.east) -- ++(0.45,0) |- node[lbl, pos=0.22] {0..1 : *} ([yshift=-0.35cm]visit.west);
\draw[relopt] (admin.east) -- ++(0.65,0) |- node[lbl, pos=0.2] {0..1 : *} ([yshift=-0.85cm]visit.west);
\draw[relopt] (admin.east) -- ++(0.65,0) |- (vdet.west);

\draw[rel] (visit.north) -- node[lbl, right=2pt] {1 : 0..*} (cond.south);
\draw[relopt] (visit.south) -- node[lbl, right=2pt] {1 : 0..*} (vdet.north);
\draw[rel] ([yshift=-0.7cm]visit.west) -- ++(-0.32,0) |- node[lbl, pos=0.72] {1 : 0..*} (proc.west);

\draw[rel] (cond.east) -- ++(0.45,0) |- node[lbl, pos=0.62, above=1pt] {source} ([yshift=0.45cm]concept.west);
\draw[relopt] ([yshift=-0.75cm]cond.east) -- ++(0.35,0) |- node[lbl, pos=0.72, above=1pt] {std} (crel.west);
\draw[relopt] (crel.north) -- node[lbl, right=2pt] {Maps to} (concept.south);
\draw[rel] (proc.east) -- ++(0.75,0) |- node[lbl, pos=0.14] {PCS / OPCS-4} ([yshift=-0.55cm]concept.west);

\draw[flow] (cdmbox.east |- adapter.west) -- node[lbl, above=2pt] {scoped stays} (adapter.west);
\draw[flow] ([yshift=0.95cm]adapter.east) -- ++(0.85,0) |- node[lbl, pos=0.72, above=1pt] {\texttt{msdrg\_group()}} (out_msdrg.west);
\draw[flow] ([yshift=0.0cm]adapter.east) -- ++(0.85,0) |- node[lbl, align=center, pos=0.72, above=1pt] {\texttt{hrg\_group\_apc()}\\ \texttt{/ pso\_flag\_apc()}} (out_hrg.west);
\draw[flow] ([yshift=-0.95cm]adapter.east) -- ++(0.85,0) |- node[lbl, pos=0.72, above=1pt] {quality audit} (out_qual.west);

\end{tikzpicture}}
\caption{\textbf{OMOP CDM tables used by \texttt{opendrg} and the resulting casemix outputs.} Solid boxes inside the dashed database boundary denote core OMOP tables used in every run; dashed boxes denote optional tables (study cohorts, mortality records, multi-episode consultant stays for HRG4+, provider/site details, and vocabulary fallback mappings). Both MS-DRG and HRG4+/PSO outputs carry an attached data-quality report.}
\label{fig:omop_er}
\end{figure*}

{\sloppy Because real-world hospital OMOP databases often omit billing fields, \textbf{opendrg} never guesses silently; every assumption is explicit and audited in an attached quality report (\texttt{omop\_\allowbreak msdrg\_\allowbreak quality()} / \texttt{omop\_\allowbreak hrg\_\allowbreak quality()}):
\begin{enumerate}[leftmargin=*,nosep]
    \item \textbf{Principal diagnosis strategy.} By default (\texttt{"status"}), only stays with an OMOP primary-diagnosis flag are grouped; stays without one are marked ungroupable. Analysts can explicitly choose a study disease list (\texttt{"codelist"}), inherit each cohort's defining condition from \texttt{CohortConstructor} (\texttt{"cohort\_\allowbreak codelist"}), test every diagnosis in turn as a sensitivity check (\texttt{"all"}), or opt into first-code ordering (\texttt{"first"}).
    \item \textbf{Code quality and vocabulary fallback.} Source ICD-10-CM / ICD-10 concepts are preferred, excluding free-text NLP rows by default so unverified notes do not inflate severity. An optional \texttt{vocabulary\_\allowbreak fallback = TRUE} path translates standard OMOP concepts when source codes are absent. Non-billable category headers (e.g.\ \texttt{N18}) are reported separately from invalid codes.
    \item \textbf{Procedures and discharge disposition.} Both standard and source procedure concepts are read, flagging stays with unmapped local procedures. Discharge status follows four tiers: in-stay death in \texttt{death}, optional per-site lookup map, standard OMOP discharge concepts, and routine home fallback.
\end{enumerate}
\par}

\subsection{Real-world OMOP portability studies (Spain and England)}

To test real-world portability across health systems, we structured a parallel two-country evaluation on live OMOP CDM databases:

\textbf{Spain (completed): MS-DRG on a hospital OMOP CDM.} We evaluated \textbf{opendrg} read-only on a live OMOP CDM v5.4 database at a Spanish specialist urology--nephrology hospital (222{,}360 historical inpatient stays; 5{,}126 stays across 3{,}791 patients in 2024; 35{,}312 matched stays with official 3M APR-GRD v35/v36 relative weights from 2018 through August 2024; aggregate summaries only). We conducted two analyses: (i)~on all 5{,}126 stays of 2024, we compared four strategies---the initial CDM audit (\texttt{"status"}), a nine-prefix specialist codelist (\texttt{"codelist"}), a 1{,}000-stay sensitivity spread (\texttt{"all"}), and official-grade grouping with complete OMOP mapping (\texttt{"status"}); and (ii)~across all seven years (2018--2024, $n = 35{,}312$ stays), we benchmarked \textbf{opendrg} MS-DRG v44.0 against the hospital's billed 3M APR-GRD v35/v36 using four plain-English metrics:
\begin{itemize}[leftmargin=*,nosep]
    \item \textbf{Cost-weight agreement (5-fold CV $R^2$):} the percentage of stay-to-stay variation in $\log(\text{3M APR-GRD weight})$ explained by \textbf{opendrg}'s MS-DRG on unseen held-out stays.
    \item \textbf{Rank agreement (Spearman $\rho$):} how consistently \textbf{opendrg} and 3M rank costlier stays above cheaper stays (0 to 1).
    \item \textbf{Hospital Case-Mix Index (CMI) difference:} percentage difference between predicted and actual mean annual cost weight per stay (both 5-fold cross-validated in-year and prospectively calibrating on year $t-1$ to predict year $t$).
    \item \textbf{Length-of-stay variation explained and clinical crosswalk:} $R^2$ on $\log(\text{LOS}+1)$ for \textbf{opendrg} MS-DRG versus base and full 4-tier 3M APR-GRD, plus direct patient overlap across the ten largest clinical groups in 2024.
\end{itemize}

\textbf{England (planned confirmation): HRG4+ on CPRD Aurum--HES OMOP.} In collaboration with Oxford HEOR-NDORMS, a parallel read-only confirmation is planned on linked CPRD Aurum--HES data mapped to OMOP CDM (retrospective fragility-fracture cohort, $n \approx 23{,}000$ postmenopausal women, 2010--2018). The study will run \texttt{opendrg::omop\_hrg()} directly inside Oxford's secure environment---covering admitted patient care episodes and spells (\texttt{visit\_detail} FCEs) as well as outpatient consultations, under both OMOP-standard and OMOP-source strategies---and benchmark episode/spell HRGs, unbundled HRGs, and National Cost Collection unit costs directly against NHS England's official Windows HRG4+ Grouper outputs.

%----------------------------------------------------------------------------------------
%	4. RESULTS
%----------------------------------------------------------------------------------------

\section{Results}
\label{sec:results}

\subsection{Verification of CMS MS-DRG}

Table~\ref{tab:results} summarises the comparison against CMS's official software on all 9{,}864 \texttt{TESTDB} stays. Across all three bundled versions (v43.0, v43.1, and v44.0), \textbf{opendrg} achieved \textbf{100\% agreement on every output field}: final DRG, initial DRG before HAC adjustments, MDC, medical/surgical class, severity tier, and return code. For v44.0 this spans 9{,}720 grouped stays across 759 distinct DRGs, 48 external-cause rejections, 96 invalid principal diagnoses, 101 stays invalid as principal discharge diagnosis (DRG~998), and all 73 stays whose final DRG shifts due to a hospital-acquired condition---including CMS's alternative rejection (return code 10) when hospital POA-reporting status is unknown. Agreement for v43.0 and v43.1 was obtained on the first run after importing their official tables with zero engine changes. Throughput on a standard laptop is approximately 1{,}200 stays per second.

\begin{table}[!t]\centering
\caption{Agreement of \textbf{opendrg} with CMS's official grouper on all 9{,}864 TESTDB stays, per bundled version.}
\label{tab:results}
\scriptsize
\setlength{\tabcolsep}{3.0pt}
\begin{tabular}{@{}lccc@{}}
\toprule
Output & v43.0 & v43.1 & v44.0 \\
\midrule
Final DRG & 9{,}864 / 9{,}864 & 9{,}864 / 9{,}864 & 9{,}864 / 9{,}864 \\
Initial DRG (before HAC) & 9{,}864 / 9{,}864 & 9{,}864 / 9{,}864 & 9{,}864 / 9{,}864 \\
MDC & 100\% & 100\% & 100\% \\
Medical / surgical class & 100\% & 100\% & 100\% \\
Severity (MCC / CC / none) & 100\% & 100\% & 100\% \\
Return code & 100\% & 100\% & 100\% \\
\midrule
Stays changed by the HAC check & \multicolumn{3}{c}{73 (all matched)} \\
Unknown hospital POA status (RC 10) & \multicolumn{3}{c}{73 (all matched)} \\
\bottomrule
\end{tabular}
\end{table}

\subsection{Verification of NHS England HRG4+ and Prescribed Specialised Services}

Table~\ref{tab:hrgresults} presents the parallel verification for NHS England HRG4+. Across all seven care settings and all 2{,}250 official test records, \textbf{opendrg} achieved \textbf{100\% agreement on every published output field} for both the 2024/25 National Costs release and the 2025/26 Local Payment Grouper (NHS England postponed the 2026/27 release). Throughput is approximately 100 episodes per second.

\begin{table}[!t]\centering
\caption{Agreement of \textbf{opendrg} with NHS England's expected results on the official HRG4+ test data (all seven datasets), per bundled version.}
\label{tab:hrgresults}
\scriptsize
\setlength{\tabcolsep}{3.0pt}
\begin{tabular}{@{}lcc@{}}
\toprule
Dataset / Output & 2024/25 NC & 2025/26 LP \\
\midrule
Admitted patient care (FCE HRG) & 750 / 750 & 750 / 750 \\
\quad Method, dom.\ proc., unbundled, PBC, dur., quality & 100\% & 100\% \\
Admitted patient care (spell HRG) & 477 / 477 & 477 / 477 \\
\quad Method, dom.\ proc., diag., LOS, count, report flag & 100\% & 100\% \\
\midrule
Non-admitted consultations (HRG, method, dom., unb.) & 200 / 200 & 200 / 200 \\
Emergency medicine (EM HRG) & 300 / 300 & 300 / 300 \\
Adult critical care (HRG, CC days, warning flag) & 250 / 250 & 250 / 250 \\
Paediatric critical care (HRG per day) & 250 / 250 & 250 / 250 \\
Neonatal critical care (HRG per day) & 300 / 300 & 300 / 300 \\
Renal dialysis (NRD HRG) & 200 / 200 & 200 / 200 \\
\midrule
All 7 datasets (every published field) & \multicolumn{2}{c}{2{,}250 / 2{,}250 (100\%)} \\
\bottomrule
\end{tabular}
\end{table}

For Prescribed Specialised Services (PSO), verifying against official outputs uncovered two subtle discrepancies in NHS England's published workbook: one sheet (\texttt{List Members}) retained stale 2024/25 entries while \texttt{Code Sets} held current rules, and the generic provider bypass code \texttt{PSS} described in the manual was omitted from the spreadsheet lists. Accounting for both, \textbf{opendrg} achieved \textbf{100\% agreement} across all 300 inpatient episodes (including all 44 multi-episode spells and 30 validation messages) and all 200 outpatient attendances, and importing the May 2026 \emph{Version 2} update reproduced the exact published change list (one new outpatient service line, one added procedure, four provider updates). Over 1{,}100 automated checks continuously verify all three groupers and both OMOP adapters across Windows, macOS, and Linux.

\subsection{Real-world portability — Spain: MS-DRG on hospital OMOP CDM (2018--2024)}

\textbf{Readiness audit and principal-diagnosis strategies (2024; Table~\ref{tab:omop}).} In the hospital's initial OMOP extract (5{,}126 stays in 2024), 95.8\% of stays carried ICD-10-CM diagnoses (median~4, mean~5.0 codes; 24{,}257 valid billable codes and 174 category headers), but primary-diagnosis flags were unpopulated (0.0\%) and 84.3\% of stays carried unmapped local procedures (only 6.1\% mapped to ICD-10-PCS). Running a sensitivity check across every diagnosis (\texttt{"all"}, 1{,}000 stays, 4{,}913 candidate runs) showed that only \textbf{9.6\% of stays receive the same DRG regardless of which diagnosis is principal}---confirming why silent heuristics fail. Using nine specialist urology/nephrology prefixes (\texttt{"codelist"}) grouped 2{,}782 / 2{,}785 matched stays (99.9\%) in 10.5 seconds with a clean length-of-stay gradient (median 2, 4, and 5 days across NCC, CC, and MCC), though surgical DRGs were capped at 2.7\% due to unmapped procedures. Once complete OMOP mapping was in place (primary/secondary flags on 99.8\% of stays, ICD-10-PCS on 85.1\%, clean billable CIE-10-ES codes without local suffixes, and standard discharge concepts), \texttt{omop\_msdrg(..., principal = "status")} grouped \textbf{5{,}114 / 5{,}126 stays (99.8\%)} in 11.8 seconds with zero heuristics and lifted surgical DRGs to \textbf{57.0\%} (2{,}915 surgical, 2{,}199 medical).

\begin{table*}[!t]\centering
\caption{End-to-end MS-DRG v44.0 evaluation on a live hospital OMOP CDM v5.4 (Spanish urology--nephrology hospital, 2024 inpatient stays; read-only, aggregate results only).}
\label{tab:omop}
\scriptsize
\setlength{\tabcolsep}{4.5pt}
\begin{tabular}{@{}p{4.6cm}rrrccp{4.2cm}@{}}
\toprule
Strategy & Stays & Claims & Grouped (RC~0) & Surgical & Time & Key finding \\
\midrule
\texttt{status} (initial CDM audit) & 5{,}126 & 5{,}126 & 0 (0.0\%) & --- & 13.2\,s & 0\% OMOP primary flags; 95.8\% ICD-10-CM coded; 84.3\% unmapped procedures \\
\texttt{codelist} (9 urology/nephrology prefixes) & 5{,}126 & 2{,}785 & 2{,}782 (99.9\%) & 2.7\% & 10.5\,s & Clinically coherent top DRGs; median LOS NCC 2\,d, CC 4\,d, MCC 5\,d \\
\texttt{all} (sensitivity, \texttt{n\_max = 1000}) & 1{,}000 & 4{,}913 & 4{,}781 (97.3\%) & 1.5\% & 12.2\,s & 954/1{,}000 groupable; only 9.6\% DRG-invariant (12.2\% same MDC) \\
\texttt{status} (complete OMOP mapping) & 5{,}126 & 5{,}126 & \textbf{5{,}114 (99.8\%)} & \textbf{57.0\%} & 11.8\,s & 99.8\% primary flags, 85.1\% ICD-10-PCS; $R^2=0.809$ vs.\ official APR-GRD weight \\
\bottomrule
\end{tabular}
\end{table*}

\textbf{Seven-year concordance against 3M APR-GRD (2018--2024; Table~\ref{tab:omop_apr}).} Across all 35{,}312 matched inpatient stays (grouped in 74.2 seconds), \textbf{99.8\%--100.0\%} grouped cleanly every year (53\%--59\% surgical); only 48 stays (0.14\%) failed due to retired CIE-10-ES codes. On unseen held-out stays, \textbf{opendrg}'s MS-DRG captured \textbf{72\%--81\%} of the variation in official 3M APR-GRD cost weights (75\%--85\% in-sample, rising to 84\% in 2024 when adding \textbf{opendrg}'s severity flag and secondary complication counts), ranked stays with Spearman $\rho = 0.78\text{--}0.86$, and matched each year's hospital Case-Mix Index within \textbf{$\pm 0.2\%$}. Applied prospectively (calibrating weights on year $t-1$ to predict year $t$), hospital CMI stayed within \textbf{$-$0.32\% to $+$1.32\%} across stable APR-GRD v36 years (2020--2023). On length of stay, \textbf{opendrg}'s 3-tier MS-DRG explained \textbf{32\%--38\%} of $\log(\text{LOS}+1)$ variation---exceeding base APR-GRD (27\%--32\%) and capturing 80\%--90\% of full 4-tier 3M APR-GRD (36\%--43\%).

\begin{table*}[!t]\centering
\caption{Seven-year portability benchmark: \textbf{opendrg} MS-DRG v44.0 (\texttt{principal = "status"}) versus the hospital's official 3M APR-GRD v35/v36 cost weights and length of stay ($n = 35{,}312$ stays, 2018--2024; positive billed APR weights). Cost-weight agreement is 5-fold out-of-sample $R^2$ on $\log(\text{APR weight})$; hospital CMI diff is the 5-fold cross-validated difference in mean hospital Case-Mix Index.}
\label{tab:omop_apr}
\scriptsize
\setlength{\tabcolsep}{5.0pt}
\begin{tabular}{@{}lrrcccc@{}}
\toprule
 & & & \multicolumn{3}{c}{Agreement with 3M APR-GRD cost weight} & \multicolumn{1}{c}{LOS variation explained} \\
\cmidrule(lr){4-6}\cmidrule(l){7-7}
Year & Stays & Grouped & Cost-weight agreement (CV $R^2$) & Rank agreement ($\rho$) & Hospital CMI diff & \textbf{opendrg} vs.\ full 3M APR \\
\midrule
2018 & 5{,}622 & 99.9\% & 72\% & 0.78 & $+$0.01\% & 33\% vs.\ 39\% \\
2019 & 5{,}716 & 99.9\% & 78\% & 0.85 & $-$0.19\% & 32\% vs.\ 40\% \\
2020 & 4{,}771 & 99.8\% & 81\% & 0.86 & $-$0.10\% & 38\% vs.\ 41\% \\
2021 & 5{,}296 & 99.8\% & 76\% & 0.84 & $-$0.06\% & 33\% vs.\ 38\% \\
2022 & 5{,}175 & 99.9\% & 80\% & 0.85 & $+$0.05\% & 33\% vs.\ 36\% \\
2023 & 5{,}311 & 99.9\% & 80\% & 0.85 & $+$0.06\% & 33\% vs.\ 38\% \\
2024\textsuperscript{a} & 3{,}421 & 100.0\% & 76\% & 0.83 & $-$0.09\% & 34\% vs.\ 43\% \\
\midrule
\multicolumn{7}{@{}l@{}}{\textsuperscript{a}January--August 2024 (latest official APR-GRD billing extract). Prospective year-over-year CMI diff (2020--2023): $-$0.32\% to $+$1.32\%.} \\
\bottomrule
\end{tabular}
\end{table*}

\textbf{Clinical crosswalk (Table~\ref{tab:omop_xwalk}).} Inspecting the ten largest inpatient groups in 2024 side by side confirms that MS-DRG and 3M APR-GRD place \textbf{90\%--100\% of the exact same patients} into matching clinical categories across every major surgical and medical procedure (e.g.\ both MS-DRG~714 and MS-DRG~713 map 100.0\% into \texttt{APR36-482}, cleanly separating uncomplicated stays at mean weight 0.649 from complicated stays at 0.754). Divergence is limited to broad ``other/miscellaneous'' medical catch-all groups (\texttt{699}/\texttt{700}), which CMS and 3M partition along slightly different boundaries.

\begin{table*}[!t]\centering
\caption{Clinical crosswalk of the ten largest inpatient groups on the Spanish hospital OMOP CDM (Jan--Aug 2024, $n = 3{,}360$ grouped stays): \textbf{opendrg} MS-DRG v44.0 versus the hospital's official 3M APR-GRD v36.}
\label{tab:omop_xwalk}
\scriptsize
\setlength{\tabcolsep}{4.5pt}
\begin{tabular}{@{}p{7.4cm}crccr@{}}
\toprule
\textbf{opendrg} MS-DRG (code \& title) & Class & Stays & Dominant 3M APR-GRD & Share in 3M group & Mean 3M wt \\
\midrule
\textbf{661} — Kidney \& ureter procedures for non-neoplasm w/o CC/MCC & Surgical & 285 & \texttt{APR36-443} & 93.3\% & 0.932 \\
\textbf{699} — Other kidney \& urinary tract diagnoses w/ CC & Medical & 187 & \texttt{APR36-466} & 50.3\% & 0.640 \\
\textbf{670} — Transurethral procedures w/o CC/MCC & Surgical & 174 & \texttt{APR36-446} & 93.1\% & 0.682 \\
\textbf{660} — Kidney \& ureter procedures for non-neoplasm w/ CC & Surgical & 171 & \texttt{APR36-443} & 95.9\% & 1.018 \\
\textbf{694} — Urinary stones w/o MCC & Medical & 132 & \texttt{APR36-465} & 92.4\% & 0.553 \\
\textbf{690} — Kidney \& urinary tract infections w/o MCC & Medical & 125 & \texttt{APR36-463} & 90.4\% & 0.468 \\
\textbf{714} — Transurethral prostatectomy w/o CC/MCC & Surgical & 115 & \texttt{APR36-482} & \textbf{100.0\%} & 0.649 \\
\textbf{713} — Transurethral prostatectomy w/ CC/MCC & Surgical & 106 & \texttt{APR36-482} & \textbf{100.0\%} & 0.754 \\
\textbf{700} — Other kidney \& urinary tract diagnoses w/o CC/MCC & Medical & 104 & \texttt{APR36-468} & 69.2\% & 0.467 \\
\textbf{669} — Transurethral procedures w/ CC & Surgical & 95 & \texttt{APR36-446} & 87.4\% & 0.787 \\
\bottomrule
\end{tabular}
\end{table*}

\subsection{Real-world portability — England: HRG4+ on CPRD--HES OMOP CDM (Pending)}

\emph{[Pending confirmation by Oxford HEOR-NDORMS.]} Live UK OMOP confirmation on linked CPRD Aurum--HES data is scheduled as the next collaborative milestone. This subsection will report aggregate results generated inside Oxford's secure environment without moving patient-level data:
\begin{itemize}[leftmargin=*,nosep]
    \item \textbf{Grouping completeness and quality audit:} share of inpatient episodes, spells, and outpatient attendances grouped cleanly (non-\texttt{UZ01Z}), ICD-10/OPCS-4 multi-domain capture, and specialty coverage under \texttt{omop\_hrg\_quality()}.
    \item \textbf{Agreement with NHS England's Windows HRG4+ Grouper:} record-by-record concordance on core episode HRG, spell HRG, dominant procedure, and unbundled HRGs under both OMOP-standard and OMOP-source strategies.
    \item \textbf{Downstream cost concordance:} mean two-year hospital cost difference using \textbf{opendrg} HRG4+ paired with NHS National Cost Collection tariffs compared against source CPRD--HES estimates.
\end{itemize}

%----------------------------------------------------------------------------------------
%	5. DISCUSSION
%----------------------------------------------------------------------------------------

\section{Discussion}

\textbf{Portability across platforms, databases, and countries.} \textbf{opendrg} runs anywhere R runs---Windows, macOS, Linux, hospital servers, or secure cloud research environments---with no Windows desktop, Java runtime, or proprietary software required. Each classification version is under 300~KB, allowing multiple fiscal years to coexist in one study. For English costing and HEOR teams, it provides the first open HRG4+ and PSO implementation that runs outside Windows; for federated OMOP networks, it brings grouping directly to the database without exporting patient records; and for countries like Spain that code in CIE-10-ES (ICD-10-CM/PCS) while reimbursing under proprietary APR-GRD, it offers an open, auditable research casemix (alongside a dedicated Spanish RAE-CMBD register adapter, \texttt{cmbd\_to\_msdrg()}, and synthetic 1{,}000-contact discharge dataset).

\textbf{Practical scope and limitations.} Four boundaries should be kept clear: (1)~\textbf{opendrg} assigns casemix groups, not tariffs or unit prices (which live in national schedules such as NHS National Cost Collection or downstream HEOR packages); (2)~MS-DRG does not replace statutory 3M APR-GRD or ambulatory IR-GRD for Spanish public reimbursement, and does not include Medicare Code Editor (MCE) coding edits or a separate Risk of Mortality score; (3)~a handful of rare paths not exercised by CMS or NHS England's test files (e.g.\ sex-diagnosis conflicts, paediatric critical-care isolation lists) are documented as open coverage items; and (4)~current HRG4+ throughput ($\sim$100 episodes/s) is well suited to hospital cohorts, while vectorisation and live UK OMOP validation (e.g.\ CPRD--HES) are the next priorities.

\textbf{Open licensing and reproducibility.} Package code is MIT-licensed. CMS tables are US federal public-domain works; NHS England HRG4+ and PSO tables are published under the Open Government Licence v3.0~\cite{ogl3}. Official classification descriptions are not redistributed. All importers, verification harnesses, specifications, and test fixtures are openly versioned so any analyst or regulator can audit results and import new annual releases.

%----------------------------------------------------------------------------------------
%	6. CONCLUSION
%----------------------------------------------------------------------------------------

\section{Conclusion}

Hospital casemix grouping is deterministic, financially consequential, and fully defined by rule tables that CMS and NHS England already publish. Treating those official tables as executable data and verifying every output record-by-record against the reference tools, \textbf{opendrg} achieves 100\% agreement for CMS MS-DRG across three versions (9{,}864 stays), 100\% agreement for NHS England HRG4+ across all seven care settings (2{,}250 records) and two versions, and 100\% agreement for Prescribed Specialised Services (500 records). Demonstrated across seven years and 35{,}312 inpatient stays in a live Spanish hospital's OMOP CDM, it proves that open, table-driven casemix grouping is genuinely portable across platforms, federated databases, and health systems---turning a closed black box into a reproducible, inspectable part of the analytical record.

%----------------------------------------------------------------------------------------
%	REFERENCE LIST
%----------------------------------------------------------------------------------------

\bibliographystyle{unsrt}
\bibliography{references}

\end{document}